% Full editable text transcription extracted from the supplied PDF.
% This preserves the wording and page-by-page reading order. Mathematical
% expressions remain in their extracted visual form rather than being rebuilt
% as semantic amsmath equations. Compile with XeLaTeX.
\documentclass[letterpaper]{article}
\usepackage[margin=0.30in]{geometry}
\usepackage{fontspec}
\usepackage{fvextra}
\setmonofont{FreeMono}
\pagestyle{empty}
\setlength{\parindent}{0pt}
\begin{document}
% ---- Original PDF page 1 ----
\begin{Verbatim}[fontsize=\fontsize{6.35}{7.15}\selectfont,breaklines=false]
 A Step-Size Resolution Closure for IMEX Integration
                          of
   Two-Dimensional Quasi-Geostrophic Turbulence:
           Truncation-Error Analysis, Wiener-Conditioned Learned Stencils,
                        and Von Neumann–Certified Stability

                     Sanaa Mouzahir1 and [TODO: co-authors / advisor]1

    1
        Department of Mechanical Engineering and Center for Computational Science and Engineering„
                       Massachusetts Institute of Technology, Cambridge, MA 02139


                                                July 14, 2026


                                                  Abstract
        We treat the time-discretization error of a fixed numerical scheme as a closure problem. For the
        spectral Adams–Bashforth-2 / Crank–Nicolson (AB2–CN2) IMEX integrator applied to two-
        dimensional quasi-geostrophic (QG) turbulence, we derive the local truncation error in closed
        form to O(∆T 6 ) and compare it, term by term, with those of AB4–CN2 and classical RK4.
        Each truncation operator splits into an analytical part (powers of the Fourier-diagonal linear
        operator, free at inference) and an expensive part built from time derivatives of the advective
        Jacobian; a compact (≈3,700-parameter) physics-initialized network estimates only the latter
        from a seven-level time stencil. The estimator’s error is analyzed exactly: the finite-difference-
        in-time error at every order is driven by the single field ω (S) , and a Wiener (least-squares)
        analysis shows the optimal learned stencil correction varies across step sizes and flow regimes
        as (∆T σ(κ))S−k , where σ(κ) is a per-shell decorrelation rate measurable from the model’s
        own inputs to O(10−3 ) relative accuracy. This yields a conditioned architecture whose pooled
        training floor deletes the dominant (pooled-variance) term of the unconditioned floor. Stability
        is handled symmetrically: the modified-equation analysis shows that a closure targeting the
        exact flow has zero linear stability margin on the advection axis, while targeting fine-step RK4
        buys a provable y 6 /144 margin; a differentiable frozen-coefficient von Neumann certificate of the
        closed scheme is derived and used both as an a priori diagnostic and as a training penalty. The
        analytic (network-free) R3 closure is unconditionally stable in our experiments and improves
        coarse-step accuracy by one to two orders of magnitude; the learned closure attains a replicated
        median 16× error reduction over ten initial conditions at ∆T = 5 × 10−3 on a 16-step horizon,
        while a horizon study shows the gain is not yet unconditional — blow-ups reappear at ∼2×
        the fine-tuned horizon — which we diagnose (a network-noise history-contamination loop) and
        attack with the certificate penalty. We report the complete error, cost, and stability budget of
        the method, including its failure modes.
        Key words. time-integration error, modified equation, closure modeling, IMEX schemes,
        Wiener filtering, von Neumann stability, pseudo-spectral methods, quasi-geostrophic turbulence,
        scientific machine learning

1. Introduction. Numerical simulations of geophysical turbulence are limited by the time step.
In stiff pseudo-spectral solvers the per-step cost is dominated by the nonlinear (advective) term,

                                                       1
\end{Verbatim}
\newpage
% ---- Original PDF page 2 ----
\begin{Verbatim}[fontsize=\fontsize{6.35}{7.15}\selectfont,breaklines=false]
and the total cost scales as T /∆t to reach a horizon T . For an integrator whose accuracy—not
linear stability—sets ∆t, halving the usable step doubles the cost. Reducing this time-discretization
cost without changing the discretization itself is the problem addressed here.
    Data-driven closures are by now standard for unresolved spatial scales on coarse grids, in engi-
neering large-eddy simulation and geophysical modeling alike [8, 9, 13]. The analogous problem for
unresolved temporal scales—the error incurred by taking a large time step with a fixed scheme—has
received far less attention. We treat it as a closure problem: derive the scheme’s local truncation
error (LTE) in closed form, identify the single piece that is expensive to evaluate, and learn only
that piece.

Goal.   The target of the construction is stated once and drives every design choice in the paper:

      run at the per-step cost of AB2–CN2 with a coarse step ∆T , while recovering the
      accuracy—and, deliberately, the stability margin—of fine-step RK4 at hf = ∆T /K.

The scheme and its formal order are unchanged; a learned correction makes one coarse step behave
as if the same trajectory had been advanced by K fine RK4 steps. Because a fine-step RK4
trajectory coincides with the exact flow to near machine precision, the K → ∞ limit targets the
exact flow; Section 5 shows why the RK4 target, not the exact flow, is the right deployment choice:
on the advection axis the exact flow has zero stability margin, while RK4’s faint dissipation buys
a provable y 6 /144 buffer against both the analytic truncation tail and the network error.

Relation to classical ideas. The construction sits between three classical methods while re-
maining distinct from each. Modified-equation analysis [17] supplies the truncation machinery but
is normally used to characterize or add order, not to emulate a finer-step version of the same scheme.
Defect correction [14] iteratively matches a higher-order scheme at the same step; we match the
same scheme (or RK4) at a smaller step, a target with no defect-correction analogue—the gain is
set by K at training time, not bounded by a reference order. Parareal [12] attains fine-propagator
accuracy at coarse cost by iterating over the horizon; our closure is an amortized parareal in which
the fine propagator is paid for once, during training, and reused at inference without iteration.
High-order time integration that pays once for the curvature of a constrained problem was devel-
oped for dynamical low-rank systems in the same group [6]. On the learning side, the estimator
belongs to the learned-discretization lineage [2, 11]: a small network constrained to the algebraic
structure of the operator it replaces, not a generic operator learner.

Contributions. (i) Closed-form LTEs of AB2–CN2, AB4–CN2 and RK4 for the QG vorticity
equation, in a shared elementary-differential language, verified symbolically and by convergence
studies (Sec. 4.1). (ii) A quantitative validity theory for any such closure: the truncation series has
a finite convergence radius ∆T ⋆ , measured per regime and collapsed to a universal law ∆T ⋆ = C τλ ,
C = 2.08 ± 0.15, with τλ the Taylor microscale of the nonlinear term; a finite-lag stencil adds an
inner wall at span ≈ ∆T ⋆ ; and error propagation through the closure is benign (1:1 in the N̈ error)
across the whole trained envelope (Sec. 4.2). (iii) A cost and stability budget of the schemes and
their closed variants, including exact amplification factors for the IMEX family and the margin
asymmetry between the exact-flow and RK4 targets (Secs. 4.3–4.4, 5). (iv) A structure-matched
estimator of the expensive LTE part (3,700 parameters, exact-physics initialization) together with
an exact error analysis of the estimator itself, culminating in a Wiener floor decomposition that
motivates—and quantifies the gain of—physics conditioning of the learned stencils on a per-sample
measurable rate σ̂(κ) (Sec. 6). (v) A differentiable frozen-coefficient von Neumann certificate for the


                                                  2
\end{Verbatim}
\newpage
% ---- Original PDF page 3 ----
\begin{Verbatim}[fontsize=\fontsize{6.35}{7.15}\selectfont,breaklines=false]
closed scheme, used as a training penalty (Sec. 7). (vi) An honest experimental record: replicated
16× coarse-step error reduction at ∆T = 5 × 10−3 over 16-step horizons, an analytic-closure ceiling
of 19–130×, and the two failure modes that remain—the large-∆T stability gap and the horizon
wall—with their diagnosed mechanism (Sec. 9).

Notation.     Table 1 collects recurring symbols.

                                          Table 1: Key notation.
        symbol         meaning
        ω, ψ           vorticity and streamfunction, ψ = ∆−1 ω
        L              stiff linear operator ν∆ − μ − β∂x ∆−1 (Fourier-diagonal)
        N (ω)          nonlinear term −J (ψ, ω) + F , F time-independent
        J (a, b)       advective Jacobian ∂x a ∂y b − ∂y a ∂x b
        B(a, b)        symmetric bilinear form with B(ω, ω) = −J (∆−1 ω, ω)
        Q(k) , Q̇      k-th total time derivative of Q along the trajectory
        ∆T, hf , K     coarse step, fine step, ratio K = ∆T /hf ; h a generic step
        S              time-stencil depth (number of stored levels; S = 7 in production)
        τ              per-step local truncation error; Rp , Dp its order-p operator and denominator
        ∆T ⋆ , τλ      convergence radius of the truncation series; Taylor microscale of N
        σ(κ)           per-shell decorrelation rate kω̇κ k / kωκ k; σ̂ its 2-mark estimate
        zL , z N , z   per-mode eigenvalues hλL , hλN , z = zL + zN
        g(z)           amplification factor; R(z) the RK4 stability polynomial
        fNN            learned (expensive) part of the closure



Outline. Section 2 states the governing equations and the exact derivative recursions every later
formula is built from. Section 3 defines the three schemes. Section 4 analyzes them: error (4.1–
4.2), cost (4.3), and stability (4.4). Section 5 assembles the step-size resolution closure and its two
possible targets. Section 6 presents the learned estimator: architecture (6.1), exact error analysis
and Wiener conditioning (6.2–6.3), and the loss including the von Neumann penalty (6.4). Section 7
analyzes the stability of the closed scheme. Section 8 documents the ensemble training pipeline.
Section 9 reports a priori and a posteriori results, cost, and failure modes; Section 10 concludes.

2. Governing equations.

2.1. The QG vorticity equation. We consider the two-dimensional quasi-geostrophic vorticity
equation on a doubly-periodic domain,
                 ∂t ω = L ω + N (ω),      N (ω) = −J (ψ, ω) + F,      ψ = ∆−1 ω,     Ḟ = 0,           (1)
with L = ν∆ − μ − β∂x ∆−1 collecting viscosity, linear drag, and the β-plane term, Fourier symbol
                                 L̂(k) = −νκ2 − μ + iβkx /κ2 ,      κ = |k|,                           (2)
and J (a, b) = ∂x a ∂y b − ∂y a ∂x b. Spatial discretization is pseudo-spectral with 2/3 dealiasing of
the quadratic products [3, 5]; the physical setting follows Vallis [16]. For the one-step analysis it is
convenient to write (1) as an autonomous ODE,
                                                                                           
           ω̇ = G(ω) = Lω + B(ω, ω) + F,          B(a, b) = − 21 J (∆−1 a, b) + J (∆−1 b, a) ,      (3)
with B symmetric bilinear and B(ω, ω) = −J (∆−1 ω, ω), so N = B(ω, ω) + F .

                                                     3
\end{Verbatim}
\newpage
% ---- Original PDF page 4 ----
\begin{Verbatim}[fontsize=\fontsize{6.35}{7.15}\selectfont,breaklines=false]
2.2. Exact derivative recursions. Repeated differentiation of (1) along the trajectory gives the
mutually recursive relations used throughout:

                            ω (k) = L ω (k−1) + N (k−1) ,              ψ (k) = ∆−1 ω (k) ,                           (4)
                                          m  
                                          X  m                        
                           N   (m)
                                     =−              J ψ (m−j) , ω (j) ,          m ≥ 1,                             (5)
                                                 j
                                          j=0

where the static forcing drops out of every N (m) , m ≥ 1 (bilinearity of J and the Leibniz rule).
Unrolled,
                               X
                               k−1
                ω (k) = Lk ω +     Lk−1−j N (j) ,   e.g. Ṅ = −J (ψ̇, ω) − J (ψ, ω̇).         (6)
                                 j=0

Each additional time derivative of N costs another layer of nested Jacobians; this nested-Jacobian
cost is what the closure will learn to avoid. Finally, since B is bilinear, the Fréchet expansion of G
terminates exactly:

       G(ω + p) = f + Jp + B(p, p),              f ≡ G(ω),           J ≡ G′ (ω) = L + 2B(ω, · ),         G′′′ ≡ 0.   (7)

3. Time-integration schemes.

3.1. AB2–CN2. Crank–Nicolson treats the stiff linear term implicitly and Adams–Bashforth-2
the nonlinearity explicitly [1, 7]:
                                                                            1 + h2 L               h
              ω n+1 = r(h) ω n + s(h)         2N − 2N
                                              3 n  1 n−1
                                                                 ,       r=              ,   s=            .         (8)
                                                                              1 − h2 L            1 − h2 L

All linear operators are diagonal in Fourier space, so the implicit solve is an elementwise division;
the only nonlinearity is the Jacobian, evaluated pseudo-spectrally at five FFTs per call with 2/3
dealiasing. The scheme is second-order accurate in time.

3.2. AB4–CN2. The four-level variant replaces the AB2 bracket:

                                                 55N n − 59N n−1 + 37N n−2 − 9N n−3
                    ω n+1 = r(h) ω n + s(h)                                         ,                                (9)
                                                                 24
with the same r, s. Because the CN pair limits the overall order, AB4–CN2 is still second order;
what changes is the composition of the truncation operators (Sec. 4.1.3) and the explicit stability
region (Sec. 4.4).

3.3. RK4. The classical fourth-order Runge–Kutta scheme, fully explicit on G:

k1 = G(ω n ), k2 = G(ω n + h2 k1 ), k3 = G(ω n + h2 k2 ), k4 = G(ω n +hk3 ),
                                                                           ω n+1 = ω n + h6 (k1 +2k2 +2k3 +k4 ).
                                                                                                  (10)
RK4 at a fine step hf serves two roles: the training-data generator (its trajectory is the exact flow
to ∼ 10−23 at hf = 10−3 and below) and one of the two possible closure targets.

4. Analysis of the schemes.

4.1. Error analysis.

                                                         4
\end{Verbatim}
\newpage
% ---- Original PDF page 5 ----
\begin{Verbatim}[fontsize=\fontsize{6.35}{7.15}\selectfont,breaklines=false]
4.1.1. Exact one-step map. Expanding the exact flow in total derivatives (equivalently, the Duhamel
integral in φ-functions) and substituting (6),
                                               2                       3
         ω(tn + h) = ω + h(Lω + N ) + h2 (L2 ω + LN + Ṅ ) + h6 (L3 ω + L2 N + LṄ + N̈ )
                         4                                ...
                     + h24 (L4 ω + L3 N + L2 Ṅ + LN̈ + N )
                        h5
                                                              ...
                     + 120  (L5 ω + L4 N + L3 Ṅ + L2 N̈ + L N + N (4) ) + O(h6 ).                       (11)

4.1.2. RK4. Because the Fréchet expansion (7) terminates, every RK4 stage can be expanded
exactly in the elementary differentials {J k f, B(·, ·)}; matching against the exact derivatives

                ω (3) = J 2 f + 2B(f, f ),     ω (4) = J 3 f + 2JB(f, f ) + 6B(f, Jf ),   ...            (12)

shows agreement through h4 and yields the first defect at h5 [4, 10].
Proposition 1 (LTE of RK4). For (3),
         h                                                                                                i
τRK4 = h5 120
           1
              J 4 f − 240
                       1
                          J 2 B(f, f )+ 120
                                         1               1
                                            JB(f, Jf )− 60 B(f, J 2 f )+ 120
                                                                          1                   1
                                                                             B(f, B(f, f ))− 80 B(Jf, Jf ) +O(h6 ) ≡ h5 T5 +O
                                                                                                          (13)
The J 4 f coefficient is 1/120: no stage composition builds the depth-four chain, so RK4 misses the
deepest elementary differential entirely. The full stage-by-stage derivation is in Appendix A.
                                                                                                     2
4.1.3. The IMEX pair. Taylor-expanding the AB brackets ( 23 N n − 21 N n−1 = N + h2 Ṅ − h4 N̈ +
h3
   ...
12 N − . . . ), the geometric expansions of r, s, and subtracting from (11):

Proposition 2 (LTE of AB2–CN2).
                                      h3 (2) h4 (2) h5 (2)    h6 (2)
                           τAB2 = −      R3 − R4 −      R5 −     R + O(h7 ),                             (14)
                                      12     24     240      1440 6

               (2)
             R3 = L3 ω + L2 N + LṄ − 5N̈ ,                                                              (15)
              (2)                                 ...
             R4 = 2L4 ω + 2L3 N + 2L2 Ṅ − 4LN̈ + N ,                                                    (16)
              (2)                                          ...
             R5 = 13L5 ω + 13L4 N + 13L3 Ṅ − 17L2 N̈ + 8L N − 7N (4) ,                                  (17)
              (2)                                              ...
             R6 = 43L6 ω + 43L5 N + 43L4 Ṅ − 47L3 N̈ + 28L2 N − 17LN (4) + 4N (5) .                     (18)

The orders h0 , h1 , h2 cancel identically (second-order accuracy).
Proposition 3 (LTE of AB4–CN2).
                                           h3 (4) h4 (4) h5 (4)
                                τAB4 = −      R3 − R4 −     R + O(h6 ),                                  (19)
                                           12     24     720 5
                     (4)                              (4)
                R3 = L3 ω + L2 N + LṄ ,       R4 = 2L4 ω + 2L3 N + 2L2 Ṅ + LN̈ ,                       (20)
                       (4)                                          ...
                      R5 = 39L5 ω + 39L4 N + 39L3 Ṅ + 24L2 N̈ + 9L N − 251 N (4) .                      (21)
                                              (2)    (4)
Remark 1 (Structural blind spots). R3 − R3 = −5N̈ : the AB2 stencil is structurally blind to
N̈ , and the weight −5 makes N̈ the dominant, irreducible part of the AB2–CN2 closure (quantified
in Sec. 4.2). AB4–CN2 zeroes the N̈ slot at orders 3–4 but pays at h5 with the enormous −251 N (4)
coefficient—the deeper explicit stencil defers, rather than removes, the nonlinear-derivative burden.
Both schemes remain second order: the leading L-monomials are set by the CN pair.

                                                       5
\end{Verbatim}
\newpage
% ---- Original PDF page 6 ----
\begin{Verbatim}[fontsize=\fontsize{6.35}{7.15}\selectfont,breaklines=false]
Remark 2 (Padé consistency check). The Lp ω coefficients of AB2–CN2, {− 12    1
                                                                              , − 12
                                                                                   1
                                                                                     , − 240
                                                                                          13
                                                                                             , − 1440
                                                                                                  43
                                                                                                      },
                                                                         hL
are exactly the error coefficients of the (1, 1)-Padé approximant r(h) to e (pure Crank–Nicolson
truncation). All coefficients in Props. 2–3 were verified symbolically.

Remark 3 (K-fine reference). Targeting K fine steps of the same scheme at hf = ∆T /K rather
than the exact flow multiplies the order-p term by (1 − 1/K p−2 )·(scheme factors); for the leading
term the factor is (1 − 1/K 2 ), and for K ≳ 100 all such factors are ≈ 1. Targeting fine RK4
subtracts τRK4 , which first differs from the exact flow at O(h5f ) = O(∆T 5 /K 5 )—negligible. The
exact-flow operators are therefore what one implements in all cases; the choice of target matters
only through the stability margin (Sec. 5).

4.1.4. Numerical verification. For a developed-flow snapshot ω n we form the exact history and fu-
ture by fine RK4, take one AB2–CN2 step of size h, and measure the empirical defect and the resid-
uals after subtracting the analytical prediction kept through order P . Keeping R3 , . . . , RP leaves a
residual O(hP +1 ); the measured log–log slopes are 4, 5, 6 for P = 3, 4, 5, confirming R3 , R4 , R5 (R6 ’s
slope-7 line sits at the float64 roundoff floor and is confirmed symbolically instead). Chain-rule
derivatives (5) are validated against RK4 finite differences to ≤ 0.7% (≤ 0.6% for N (5) ).




           Figure placeholder: vs-exact-flow convergence. kτemp k (slope 3) and residuals
              after subtracting R3 , R3 +R4 , R3 +R4 +R5 (slopes 4,5,6) vs h. Source:
                            analysis/validate_ab2cn2_vs_truth.py.




4.2. Magnitudes, the convergence radius, and error propagation. The truncation opera-
tors are only useful if their magnitudes are known on real turbulence. Three questions decide the
design: how fast do the N (m) grow (convergence radius); how deep may the time stencil reach
(inner wall); and how do estimation errors propagate through the assembled closure (amplifica-
tion). All measurements below: float64, 2/3-dealiased Jacobians matching the solver right-hand
side, 5122 forced-turbulence members at marks hf = 5 × 10−3 ; truth is the analytic spectral chain
rule (5); errors are relative L2 in the resolved band. Members: FRC-Re25k (ν = 4.0 × 10−5 , β = 1),
FRC-combo (ν = 4.0 × 10−5 , β = 0.5), FRC-kf4 (ν = 1.02 × 10−4 , β = 1).
4.2.1. Outer
         P wall: finite convergence radius. The defect of one AB2–CN2 step is a power series
δ(∆T ) = p≥3 cp ∆T p , cp = Rp /Dp , Dp = {12, 24, 240, 1440}. By Cauchy–Hadamard, 1/∆T ⋆ =
lim supp kcp k1/p ; the successive-ratio test only brackets ∆T ⋆ (the cp are not geometric), so the
headline number is the root-test median over p = 3, . . . , 6 (32 anchors per member).

       member         ν        β      N (m+1) / N (m) , m = 0, . . . , 4   ∆T ⋆ (root)   ratio bracket
                          −5
       Re25k     4.0 × 10      1.0   34, 58, 72, 84, 93                      0.066       [0.017, 0.136]
       combo     4.0 × 10−5    0.5   15, 27, 37, 45, 51                      0.139       [0.033, 0.270]
       kf4       1.0 × 10−4    1.0   10, 19, 28, 33, 37                      0.199       [0.043, 0.366]




                                                     6
\end{Verbatim}
\newpage
% ---- Original PDF page 7 ----
\begin{Verbatim}[fontsize=\fontsize{6.35}{7.15}\selectfont,breaklines=false]
Proposition 4 (Finite radius). ∆T ⋆ is finite and set by the growth of the N -derivatives. For
∆T > ∆T ⋆ no finite-order analytic closure converges: past ∆T ⋆ , adding correction orders makes
the error worse, regardless of how many orders are carried.

This overturns the naive expectation that higher-order corrections extend the usable step. Note
also that the per-order growth ratios increase with m (34 → 93 for Re25k): higher derivatives grow
faster than any single rate—derived next.
4.2.2. The Taylor microscale of the nonlinear term. View N (t) as a stationary process with the
domain inner product and define ρ(τ ) = EhN (t), N (t + τ )i/E kN k2 . Under stationarity, EhN, Ṅ i =
                                                                                  2
           2
2 dt E kN k = 0, and differentiating once more gives EhN, N̈ i = −E
1 d
                                                                            Ṅ        . Hence the exact identity
(no Gaussianity, no spectral assumption):

                                                                       2
                               τ2                             E Ṅ                       1
                    ρ(τ ) = 1 − σ 2 + O(τ 4 ),         σ2 =            2,        τλ =      .               (22)
                               2                              E kN k                     σ

τλ is the Taylor microscale of the N -process; the same identity holds for ω and per wavenumber
shell, with σ(κ) = kω̇κ k / kωκ k. Per-order growth ladders as the moment ratio
                                           R
                                              σ(κ)2(m+1) EN (κ) dκ
                                     σm = R
                                       2
                                                                   ,                        (23)
                                               σ(κ)2m EN (κ) dκ

which increases with m because σ(κ) increases with κ: higher derivatives progressively sample the
fast small-scale tail—exactly the escalation (34 → 93) in the table above. Verification: the identity
holds to 15% on N and 3% on ω (curvature fit vs 1/σ; Re25k τλ σN = 1.15, kf4 1.09; on ω, 0.98
and 0.97). Two further measured facts: the state decorrelates an order of magnitude more slowly
than its nonlinear term (σN /σω ≈ 7.5–11; over 27 lags ρω ≥ 0.96 while ρN falls to 0.2–0.6), so a
time stencil on the (ω, ψ)-history sees smooth data; and ρN plateaus at a positive value (≈ 0.2, the
static forcing plus time-mean of −J ), which biases integral decorrelation times but not τλ .




         Figure placeholder: measured ρN (τ ), ρω (τ ) vs the identity parabolas 1 − τ 2 σ 2 /2
                              (not fitted), Re25k and kf4. Source:
                   Theoretical_guarantees/check_decorrelation_time.py
                               (img/rho_vs_identity_*.png).




Proposition 5 (Universal radius law). ∆T ⋆ = C τλ with C = 2.08 ± 0.15, constant to 7% across
2.5× in Re and 2× in β (Re25k 2.27, combo 2.05, kf4 1.92): the analytic closure’s convergence
horizon is two Taylor microscales of the nonlinear term. All (Re, β)-dependence of ∆T ⋆ enters
through the single measurable rate σ.

4.2.3. Inner wall: the finite-lag stencil. The closure estimates N -derivatives from S stored marks
at lags {0, ∆T, . . . , (S − 1)∆T }. Depth raises the approximation order (O(∆T S−m ) for the m-th
derivative) but reaches further back in time; old marks are stale once the span approaches ∆T ⋆ .

                                                   7
\end{Verbatim}
\newpage
% ---- Original PDF page 8 ----
\begin{Verbatim}[fontsize=\fontsize{6.35}{7.15}\selectfont,breaklines=false]
Measured (N̈ , depth 4 vs depth 7, exact spectral spatial operators): at ∆T = 5 × 10−3 depth 7
helps 7.6–25× in every member; at ∆T = 1.5 × 10−2 it still helps 3–6× for combo/kf4 but hurts
for Re25k (0.99 → 1.51)—precisely where its span 6∆T = 0.09 exceeds Re25k’s ∆T ⋆ = 0.066.
Proposition 6 (Inner wall). Deeper stencils lower the N -derivative error while the stencil span
stays well inside ∆T ⋆ , and raise it once the span exceeds ∆T ⋆ ; the usable step of a depth-S closure
scales as ∆Tinner ∼ ∆T ⋆ /(S − 1). Measured at Re25k, S = 7: predicted 0.011, observed 0.013.
Consequence for pooled training: a (Re, β)-blind pool caps at the worst member’s inner wall; at
S = 7, ∆T = 1.5 × 10−2 the Re25k samples are unlearnable by any network (the field decorrelates
inside the stencil) and are retained in the pool by design as past-wall stressors.
4.2.4. Benign error propagation. An error ε on the learned field enters the closure as (∆T p /Dp )|ap,k | ε Lp−k fieldk
                                           amplified. Measured at the operating step ∆T = 5×10−3
terms carrying powers of L are potentially ...
(per-order errors 2.6/3.0/4.0% on Ṅ /N̈ / N ): the un-amplified L0 N̈ term equals kδk to three fig-
ures in every member, and the total error is 3.00%—exactly the N̈ error. At the envelope edge
∆T = 1.5 ×       −2     per-term errors, pure amplification geometry, Re25k): L0 N̈ carries 100.2%
           ... 10 (unit
of kδk, L N 2.6%, L Ṅ 0.5%, L2 Ṅ 0.00%.
         0           1

Proposition 7 (Benign amplification). Across the entire trained envelope ∆T ≤ 1.5 × 10−2 the
closure error equals the relative N̈ error, one-to-one and un-amplified. Any improvement in N̈
accuracy converts to closure-error reduction 1:1, and N̈ accuracy is the single quantity that sets the
rollout ceiling.
Two design consequences: the equal-weight per-order relative-L2 loss is the right objective across
the envelope (no amplification weighting needed), and N̈ validation error is the number to watch.
Corollary 1 (Operating envelope). For a depth-S stencil, operate with (S − 1)∆T comfortably
below ∆T ⋆ (Re, β) = 2τλ . Then the series converges (Prop. 4), the stencil resolves the derivatives
(Prop. 6), and derivative accuracy converts 1:1 (Prop. 7). At S = 7, Re25k caps the pooled envelope
near ∆T ≈ 10−2 ; combo and kf4 remain inside through 1.5 × 10−2 .

4.3. Cost analysis. Per-step cost is measured in Jacobian evaluations (five FFTs each: two inverse
transforms of ∇ψ, two of ∇ω, one forward of the dealiased product; all other operations are Fourier-
diagonal multiplies).
• AB2–CN2 / AB4–CN2: one new Jacobian per step; the AB history reuses cached N n−j .
   Implicit solve: elementwise division. Cost ≈ 5 FFTs/step.
• RK4: four Jacobian evaluations per step (≈ 20 FFTs/step), and no reuse across steps.
• Closed AB2–CN2 (Sec. 5): the bare step, the analytical R3 scaffolding (Fourier-diagonal
   multiplies on the already available ω̂, N̂ ; the L3 ω piece folds into the implicit solve at zero cost),
   and one network forward. Measured in the production stepper: bare 5 FFTs/step → closed 8
   FFTs/step, plus the network (≈ 170 multiply–accumulates per grid point—negligible next to
   an FFT). The conditioned variant adds zero further FFTs at rollout: its spectral context reuses
   the stepper’s own spectral states.
To integrate to horizon T at matched accuracy, the bare scheme needs T /hf steps at the fine step
and the closed scheme T /∆T steps at the coarse step: the ideal speedup is
                   bare FFTs/step
speedup ≈ K ×                      = K × 58             (vs fine AB2–CN2),        K × 20
                                                                                      8    (vs fine RK4),
                  closed FFTs/step
                                                                                               (24)
with K = ∆T /hf set at training time. The honest accounting of what fraction of (24) is realized—
given the measured accuracy at each ∆T —is in Sec. 9.4.

                                                    8
\end{Verbatim}
\newpage
% ---- Original PDF page 9 ----
\begin{Verbatim}[fontsize=\fontsize{6.35}{7.15}\selectfont,breaklines=false]
4.4. Stability analysis.
4.4.1. Explicit regions. For the explicit families the absolute-stability region is {z : |g(z)| ≤ 1}
                                                                                 2     3   z4
with g the max-modulus root of ρ(ζ) − zσ(ζ) (AB) or g = R(z) = 1 + z + z2 + z6 + 24           (RK4).
For advection the eigenvalues are near-imaginary, so the binding number is the imaginary-axis
stability boundary (ISB): AB2’s imaginary-axis   √ interval is empty (marginal at the origin only),
AB3’s is ≈ 0.724, AB4’s ≈ 0.43, and RK4’s is 2 2 ≈ 2.83 [4, 10]. RK4’s region is vastly larger per
step—but costs four Jacobians (Sec. 4.3); AB2 survives in the IMEX pair because the CN-treated
viscosity pulls the advective eigenvalues off the axis.




          Figure placeholder: absolute-stability regions of AB1–AB4 and RK4 with axis
                   intercepts. Source: analysis/ab_stability_regions.py.




4.4.2. IMEX amplification factors. On a Fourier mode, split z = zL + zN with zL = hL̂(k) treated
by CN and zN the explicit advection symbol; the frozen-coefficient (von Neumann) worst case over
background velocity is zN = ih U |k|. With
                                      1 + zL /2                 zN
                                 r=             ,       ρ=             ,                       (25)
                                      1 − zL /2              1 − zL /2

AB2–CN2 is the two-step recurrence ω n+1 = (r + 32 ρ) ω n − 21 ρ ω n−1 with companion amplification
                                
                  g 2 − r + 32 ρ g + 21 ρ = 0,   gAB2 = max-modulus root,                     (26)

AB4–CN2 the analogous quartic companion, and gRK4 = R(zL + zN ). Absolute stability of the
discrete mode is |g| ≤ 1; the maximum stable ∆T per regime follows by bisection over all dealiased
modes, using each regime’s measured characteristic velocity.



           Figure placeholder: per-regime stability summary—|g| = 1 contours at the
           median zL , the zL (k) cloud, worst-mode |g| vs ∆T , and max stable ∆T per
                    (regime, scheme) over the 14 ensemble regimes. Source:
                  Theoretical_guarantees/imex_stability_regions.py,
                                   analysis/vn_stability.py.




4.4.3. Truncation symbols. For the closure analysis we need the LTE in amplification space. Re-
placing Lj ω → λjL and N (m) → λN (λL + λN )m per mode, the order-p operator becomes the
polynomial
                                                 X
                                                 p
                      Pp (zL , zN ) = ap,0 zLp +   ap,k zLp−k zN (zL + zN )k−1 ,           (27)
                                                k=1


                                                    9
\end{Verbatim}
\newpage
% ---- Original PDF page 10 ----
\begin{Verbatim}[fontsize=\fontsize{6.35}{7.15}\selectfont,breaklines=false]
        (2)
with ap the AB2–CN2 coefficient rows of Prop. 2, so that
                                                        X Pp                                          z5
                 τ̂AB2 = ez − gAB2 = −                               ,       τ̂RK4 = ez − R(z) =         + O(z 6 ).          (28)
                                                        p
                                                                Dp                                   120


5. The step-size resolution closure.
                                         n+1
5.1. Definition and targets. Since ωAB2CN2        + τ = ω(tn + h) exactly, the correction to add to
each bare step is δ = τ , with two natural targets:
                           3    (2)       4       (2)       5       (2)
     δ1 = τAB2 = − h12 R3 − h24 R4 − 240
                                     h
                                         R5 − . . .                                    (target = exact flow / K → ∞),         (29)
     δ2 = τAB2 − τRK4 = δ1 − h T5             5
                                                                                       (target = fine-step RK4),              (30)
                           P
with symbols δ̂1 = −             p Pp /Dp and δ̂2 = δ̂1 − z /120.
                                                           5                           Removing the truncation through O(h6 )
leaves O(h7 ):
                                                            ∆T 3      ∆T 4      ∆T 5      ∆T 6
                        n+1
                       ωclosed    n+1
                               = ωAB2CN2 −                       R3 −      R4 −      R5 −      R6 .                          (31)
                                                             12        24       240       1440

5.2. Why target RK4: the stability margin. Closing exactly gives g = gAB2 + δ̂1 = ez + O(z 6 )
or g = gAB2 + δ̂2 = R(z) + O(z 6 ). On the advection axis z = iy,

                                                                                  y6
              |eiy | ≡ 1       (margin 0),              |R(iy)| = 1 −                + O(y 8 ) < 1    (margin > 0) :         (32)
                                                                                 144
RK4 agrees with ez through z 4 and the z 5 /120 difference is purely imaginary on the axis (a phase
error), so the modulus deficit first appears at y 6 —RK4 is faintly dissipative. The exact flow,
being neutral, buys no buffer: any perturbation of an exact-flow-matching closure (analytic tail or
network error) can be destabilizing. This asymmetry, made precise in Sec. 7, is the deployment
rationale for δ2 /fine-RK4 as the reference: the goal of the whole construction is AB2–CN2 cost at
coarse ∆T with the accuracy and stability of fine-hf RK4.

5.3. Analytical/learned split. Each operator separates into a part that is free at inference and
a part that would require nested Jacobians:

 Rp =         Rpanal           +          RpNN                  ,         e.g.   R3anal = L3 ω + L2 N,     R3NN = LṄ − 5N̈ . (33)
              | {z }                      |{z}
        powers of L on ω, N        time derivatives of N

The analytical parts are Fourier-diagonal
                                ...              (or reuse the Jacobian already computed for N ). The
learned parts contain Ṅ , N̈ , N , . . . —each another layer of nested Jacobians...(5)—and are the only
expensive pieces. The network therefore predicts the local fields [Ṅ , N̈ , N ]; every Lk -weighting
(including the nonlocal β term) is applied analytically at assembly, never learned. By Prop. 7 the
N̈ channel dominates the assembled correction 1:1 across the envelope.

5.4. Inference decomposition. At inference the linear L3 ω piece is folded into the implicit solve
(it shifts the CN denominator), the L2 N piece is applied explicitly on the cached Jacobian, and
the network supplies the rest:
                                 2                     1 2 n                
     n+1
    ωclosed = IMEX r, s ; L + ∆T    L 3
                                          + ∆T 3
                                                 1 −  1
                                                       2   L N  + ∆T 3
                                                                       1 −  1
                                                                             2   fNN (history), (34)
              |          {z    12
                                        } |          {z 12
                                                     K
                                                               } |         K
                                                                               {z            }
                       implicit, folded                             explicit, analytical             explicit, learned


                                                                          10
\end{Verbatim}
\newpage
% ---- Original PDF page 11 ----
\begin{Verbatim}[fontsize=\fontsize{6.35}{7.15}\selectfont,breaklines=false]
                                                                                           ...
with fNN ≈ 12 1
                [LṄ − 5N̈ ] the ∆T -independent learning target (plus the R4 channel via N where
warranted). The network is applied online: it is evaluated on the closure trajectory’s own evolving
history and its correction feeds the next step, so errors compound autoregressively and the network
never sees truth at deployment.

6. Method: the learned estimator.

6.1. Network architecture. The estimator (CheapDerivClosureNet) maps the stored step his-
tory to the local derivative fields,
                                                                             ...
                X = ω0 , . . . , ω−(S−1) , ψ0 , . . . , ψ−(S−1) 7−→ [Ṅ , N̈ , N ], S = 7, (35)
in four stages that mirror the operator being estimated:
 1. Time finite differences (frozen, exact). Backward-FD rows Wki from the Vandermonde
             P           p                   (k)  (k)
    system i Wki (−i)                                                                                 k
                       p! = δkp produce ωFD , ψFD , k = 0, . . . , 3, via the per-sample path Wki /∆T .
    Nothing is learned here: ... rescaling the unit stencil per sample is exact. Accuracy at S = 7: Ṅ
    O(∆T 6 ), N̈ O(∆T 5 ), N O(∆T 4 ).
 2. Order clip. Only time-orders 0, . . . , 3 are emitted; the unused rows (scaled 1/∆T 4...6 ≈ 1010 –
    1013 ) are dropped before any parameterized layer. All S snapshots still feed each emitted row
    (the order-3 row of a 7-node stencil is 4th-order accurate).
 3. Learnable spatial gradients. Depthwise width-15 stencils (14th-order central-difference ini-
    tialization) apply ∂x , ∂y to the 8 differentiated fields, circular padding. The stencil parameter is
    the dimensionless unit-spacing operator; the physical 1/dx, 1/dy scaling is applied per sample
    at forward—exact by linearity, and the optimizer never sees the 1/dx amplification. One stencil
    pair serves all grids.
 4. Jacobian features and physics-initialized mix. The (4)2 = 16 pointwise products J (ψ (i) ,ω (j) )
    feed a 1×1 mix initialized to the exact chain-rule binomials of (5) (w[m−1, (m−j, j)] = − m     j ):
    at initialization the network is the analytic estimator up to the two FD truncations.
Totals: 3,700 parameters, of which the 49 time-FD weights are an inert frozen buffer (the trained
capacity is the stencils and the mix); ≈ 170 multiply–accumulates per grid point; no FFTs inside
the model—the pointwise product is the sole irreducible nonlinearity, and carrying ψ in the history
makes even the β term local (β∂x ∆−1 ω = −βv). The architecture is exactly quadratic in the input
field, matching the quadratic nonlinearity of (3).
Remark 4 (Why the spatial stencils are learnable at all). Their job is double: approximate ∂x , ∂y
and absorb the differentiated time-FD error. The second job is the Wiener mechanism analyzed
next, and is why the width was set to 15 and why the stencils are conditioned in Sec. 6.3.

6.2. Exact error analysis of the estimator. Two error operators separate cleanly. The time-FD
estimate carries
                                                                        1 X
                                                                            S−1
                (k)
          εk ≡ ωFD − ω (k) = Ck ∆T S−k ω (S) + O(∆T S−k+1 ),       Ck =     Wki (−i)S ,         (36)
                                                                        S!
                                                                            i=0

with ε0 = 0 and εψ          −1      −1 commutes with the stencil). The learned gradients carry
                   k = ∆ εk (∆
Dθ = ∂ +δθ with δθ the stencil gap. Propagating both through the bilinear Jacobian and discarding
O(εδ, ε2 , δ 2 ),
                                                 h                           i
∆N (m) = T (m) +A(m) [δθ ],    T (m) = ∆T S−m Cm −J ∆−1 ω (S) , ω −J ψ, ω (S) +O(∆T S−m+1 ) :
                                                                                             (37)

                                                 11
\end{Verbatim}
\newpage
% ---- Original PDF page 12 ----
\begin{Verbatim}[fontsize=\fontsize{6.35}{7.15}\selectfont,breaklines=false]
the leading time error at every order m is driven by the same field ω (S) , entering through exact
Jacobians, while A(m) [δθ ] is linear in the stencil gap. Training therefore solves a Wiener problem
[18]: the least-squares-optimal stencil correction δ ⋆ pre-cancels the time error it can reach.
6.2.1. The pooled floor and its three parts. For a pool of tiers j = (member, ∆T ) with weights wj ,
orthogonal decomposition of the quadratic loss about the per-tier optima gives
                 X                         X                     2
                                                                     X                      2
      Luncond
       min    =      w j L j (δ ⋆⋆
                                j  ) + min   w j E j A[δ ⋆
                                                         j − δ̄]   +   wj Ej A[δj⋆⋆ − δj⋆ ] , (38)
                                                 δ̄
                      j                               j                                  j
                  |        {z        }       |                {z                 }   |                  {z             }
                 (i) cascade remainder           (ii) pooled-variance mismatch           (iii) width-W shape deficit

where δj⋆ is the best width-W filter for tier j and δj⋆⋆ the best shell-diagonal transfer. Measured at
S = 7, m = 2, W = 15 (valid band): diagonal-absorbable fraction 0.21–0.29; full-slot absorbable
0.36–0.44; (ii) = 0.72; (iii) = 0.31. A single unconditioned stencil can only fit the pooled mean;
its net absorption is ≈ 0.03—the pooled variance (ii) dominates everything. Deleting (ii) is what
conditioning is for.
6.2.2. Single-mode analysis of the driver. Under a sweeping split N̂ = −i(k · U)ω̂ + fˆ (A1), frozen
coefficients over the stencil span (A2), advective band    2  λ2 (A3), and subdominant triads in
                                       (n)    n
                                                    P Ωj ˆ(n−1−j)
the shell balance (A4), induction on ω̂ = a ω̂ + j a f            , a = iΩ + λ, gives the per-shell
transfer between derivative orders:

        ωd                     d
         (S) = σ (κ)S−k eiφ(κ) ω
            κ   ω
                                 (k) + r̂,
                                    κ                       φ = (S − k) π2 sign(Ω),                   r̂ = r̂1 + r̂2 + r̂3 ,   (39)

with the shell rate σω (κ)2 = hΩ2 ishell = 12 κ2 U 2 + 12 β 2 /κ2 − U β cos αU . The remainder taxonomy
bounds what any conditioning can recover: r1 (cascade input, shell-incoherent with ω (k) ; irre-
ducible—measured incoherent fraction 0.71–0.79 on healthy tiers); r2 (coefficient drift within the
stencil span, size (S−k)(S+k−1)
                         2      τκ /TU with τκ = 1/σω (κ) the per-shell state timescale and TU the
large-eddy turnover, conditionable in time); r3 (intra-shell anisotropy: for odd S − k the sweeping
transfer ∝ cosS−k θ shell-averages to zero, so an isotropic correction captures none of it and only
cos 3θ, cos 5θ harmonics are reachable by angular features).

6.3. Physics conditioning. Combining (36) and (39), the per-shell regression coefficient of the
time error on the resolved field factorizes:
                           (k)
                          rj (κ) = ∆TjS−k × Ck σω (κ; Re, β, μ)S−k eiφ(κ) +h.o.t.                                              (40)
                                   | {z } |              {z             }
                                           exact                    g (k) (κ)


Two structural facts follow. First, the optimal stencil correction varies as (∆T σ(κ))S−k —a factor
35 ≈ 243 across our ∆T -sweep at k = 2—which derives the physics conditioning: no single un-
conditioned stencil can track it. Second, (Re, β, μ) enter only through σω (κ), so the regime never
needs to be supplied as scalars: it is measurable from the model’s own input stack. The two-mark
estimator with arcsin de-biasing (a single mode rotating at Ω has FD response ∆T  2
                                                                                    | sin(Ω∆T /2)|),

                                            2                                               
                                                                   ∆T ∥(ω̂0 −ω̂−1 )κ ∥
                                 σ̂(κ) =      arcsin min            2  ∆T ∥(ω̂0 )κ ∥ , 1          ,                            (41)
                                           ∆T
recovers σ(κ) with median relative error 0.000–0.001: data-conditioning is lossless. The measured
∆T -collapse of rj /∆T S−k across the sweep is 1.6× (vs 243× raw), confirming (40).



                                                              12
\end{Verbatim}
\newpage
% ---- Original PDF page 13 ----
\begin{Verbatim}[fontsize=\fontsize{6.35}{7.15}\selectfont,breaklines=false]
Conditioned architecture (production: local taps). The deliverable keeps the zero-FFT
locality of the estimator and conditions the stencils by tap modulation:
                                       ∆T S−k                         
                     (k)      (k)
                 tapsd = tapsd,base +             ∆tapsθ,k,d σ̂-features , k ≥ 1, (42)
                                       ∆Tref

with ∆Tref = 1.5 × 10−2 anchoring the exact scaling law at O(1) (raw ∆T S−k ≈ 10−11 –10−16
renders the path numerically dead: raw ∆T S−k ≲ 5 × 10−8 down to ∼ 10−14 across the sweep),
amplitude exactly zero on the k = 0 channels (ε0 = 0: the undifferenced mark has no time error to
absorb), and ∆tapsθ a small shared MLP (10 → 24 → 24 → 240, tanh, ≈ 6.9k parameters) read-
ing [xi , log(1+xi )], xi = ∆T σ̂(κi ), at five fixed relative shells κi /κmax ∈ {0.15, 0.3, 0.5, 0.7, 0.85}.
The output head is zero-initialized: at initialization the conditioned model is bit-identical to the
                                                (k)
unconditioned control. A spectral variant D̂d = ikd [1 + ∆T S−k gθA (x, κ̃)] + kd ∆T S−k gθB (x, κ̃) has
zero shape deficit and serves as the ceiling instrument, but costs ≈ 24 FFTs per forward and is not
deployed.
                                                                                               p
Predicted ceilings. The conditioned floor deletes term (ii) of (38); per-tier, floorcond
                                                                                   j    = rawj irrj + cohj (def(F)
Predictions: kf4@1.5 × 10−2 : 0.031 → 0.023; FRC-256@1.5 × 10−2 : 0.047 → 0.037; FRC-256@10−2 :
0.0068 → 0.0055; pooled raw floor over 9 tiers 0.175 (vs the observed unconditioned plateau 0.19);
conditioned-pooled excluding the past-wall tier ≈ 0.05. Section 9.1 confronts these with the trained
result.

6.4. Loss.
                                                                                               ...
Offline derivative loss. Per-sample, per-order relative L2 with equal weights over [Ṅ , N̈ , N ]
(Prop. 7 makes equal weights correct), predictions 2/3-dealiased before the loss, and a norm floor
max(ktc k , 0.1 × member-median ktc k) in the denominator to cap the leverage of near-zero-target
samples without dropping them. Medians are logged alongside floored means: relative errors on
near-quiescent samples otherwise poison the optimizer (Sec. 8).

Rollout fine-tune. The offline objective does not see          the closed loop, so the estimator is fine-
                                                     1 PM              closed , ω truth ) against fine-RK4
tuned through the stepper: the supervised head is M       m=1 relL2 (ωm          m
truth over M coarse steps (truncated backpropagation, window 4; M scheduled up to 21), followed
by a truth-free tail. The model that produced every a posteriori result of Sec. 9 used, on that
                                           (f )  (f −1) 
tail, an annulus-enstrophy hinge relu log Zann /Zann      over the aliased band |k| > 32 kmax , weight
10−3 (larger weights destabilize the supervised head—an objective-imbalance failure we observed
directly).

Upgraded tail (current training campaign; results pending). Two replacements for the
hinge are in a running λ-sweep whose results had not landed at the time of writing and are not
claimed in Sec. 9:
 1. Analytic free-tail supervision. On truth-free rolled steps the model heads are supervised against
    the on-the-fly exact analytic targets (5) of the pre-step rolled state (weight 10−2 , capped): the
    estimator must remain a consistent derivative estimator on its own trajectory, arbitrarily far
    from the data horizon.                                                                             2
 2. Von Neumann certificate penalty (derived in Sec. 7): Lstab = λ meanκ relu |Geff (κ)| − (1 − ε) ,
    ε = 0.02, restricted to the linearization’s validity shells |∆T σ̂(κ)| ≤ 21 ; λ ∈ {0.1, 1, 10} is the
    only new hyperparameter.

                                                    13
\end{Verbatim}
\newpage
% ---- Original PDF page 14 ----
\begin{Verbatim}[fontsize=\fontsize{6.35}{7.15}\selectfont,breaklines=false]
7. Stability analysis of the learned closure.

7.1. Amplification of the closed scheme. The correction enters the update additively, hence
the amplification additively:

                           gNN,i (z) = g⋆,i (z) + ε̂i (z),    ε̂i = ε̂an,i + ε̂NN,i ,             (43)

with g⋆,1 = ez , g⋆,2 = R(z), ε̂an the deterministic leftover from cropping the closure at order p, and
ε̂NN the stochastic fit residual.

Proposition 8 (Finite-order closure: the base shows through). Assembling R3 , . . . , Rp cancels the
base scheme’s truncation series term-by-term up to order p only:

                           gNN,i = g⋆,i − τp+1
                                           base p+1
                                               z    + O(z p+2 ) + εrel O(h3 ) .                   (44)
                                        |         {z         } | {z }
                                                     ε̂an,i               ε̂NN,i


The surviving tail is the base scheme’s own (τkAB2 6= τkAB4 ; e.g. R3 carries −5N̈ for AB2 and 0
for AB4), never the target’s: above the assembled order the closure reverts to gbase , even when
matching RK4.

Corollary 2 (Empirical closure, K → ∞). The full-tail empirical correction δ = Φref − Φbase is a
difference of maps, not a truncated series: ε̂an = 0 and gNN,i = g⋆,i + εrel O(h3 ), base-agnostic.

7.2. Margin bound. By the triangle inequality, |gNN,i | ≤ 1 is guaranteed whenever εi ≤ mi (z) =
1 − |g⋆,i (z)|, tight on the advection axis. There the two targets differ qualitatively (cf. (32)):
m1 (iy) = 0—no perturbation of an exact-flow-matching closure can be certified stable—while

                                                   y     6
                                         m2 (iy) = 144 + O(y 8 ) :                                (45)

matching RK4, the closed scheme is provably stable whenever −τp+1
                                                               base (iy)p+1 +ε O(h3 ) ≤ y 6 /144.
                                                                              rel
Refining h helps the network term; the analytic-ceiling term is an h-independent comparison of
y-coefficients set by the base scheme. This margin asymmetry is the quantitative form of the
deployment rationale of Sec. 5.2.
                                                            P
7.3. Accumulated error. Over n = T /h steps, ke(T )k ≤ εi m<n |g|m , which for |g| ≤ 1 gives

match exact: ke(T )k ≤ Cbase hp T +εrel O(h2 )T ;    match RK4: ke(T )k ≤ Ch4 T +Cbase hp T +εrel O(h2 )T,
                                                                                                 (46)
(drop Cbase hp T in the empirical case). The learned floor εrel O(h2 )T preserves second-order conver-
gence, is linear in horizon and in network quality; the finite-order floor is smaller in h but frozen
(reducible only by adding analytic orders) and base-dependent.

7.4. A trainable certificate. The margin bound needs |gNN | per mode for the actual trained net-
work. Linearizing the learned piece about a frozen advecting background (the Wiener linearization:
Jacobians frozen to iσ advection, the stencil pair entering through its trigonometric-polynomial
symbol ratio ρ(κ) = Ŵbase+mod (θ)/Ŵexact (θ) evaluated on the sample’s own σ̂ shells) yields the
closed AB2–CN2 companion with effective symbols
               2                                          2                       
   I = L̂ + ∆T
             12 L̂
                   3
                     (implicit, folded),  E = iσ̂ 1 + ∆T12 L̂
                                                              2      1
                                                                  + 12 L̂ T1 − 5 T2 (explicit), (47)


                                                       14
\end{Verbatim}
\newpage
% ---- Original PDF page 15 ----
\begin{Verbatim}[fontsize=\fontsize{6.35}{7.15}\selectfont,breaklines=false]
Tc the linearized transfer of the learned Ṅ , N̈ heads, and
                                        
                                 a+b c                                                                    1
Geff (κ) = spectral radius of               ,      a = r 1+ ∆T I , b = 23 ∆T rE, c = − 21 ∆T rE, r =              .
                                  1    0                    2
                                                                                                        1 − ∆T
                                                                                                             2 I
                                                                                                     (48)
Geff is differentiable in every trainable tensor (taps, modulation, mix), reproduces the analytic-arm
stability exactly at zero modulation (ρ ≡ 1), and is evaluated only on shells with |∆T σ̂| ≤ 12 , outside
which the frozen-coefficient model provably departs from the dynamics and the penalty would be
noise. It serves as both diagnostic (a per-step contraction certificate is horizon-free, removing the
fixed-horizon asterisk of rollout training structurally) and penalty (Lstab of Sec. 6.4). Caveats: it
is linearized (the r2 , r3 remainders are neglected) and isotropic per shell. On the current stable
reference model the certificate reads |Geff | valid-shell mean 0.9992, max 0.9997—marginal at large
∆T , as it must be.

8. Ensemble training.

Data. Training data are deep 28-mark windows advanced by fine RK4 from developed snapshots
of an ensemble of forced-turbulence members spanning the (Re, β)-range (5122 , FRC-256 at 2562 ;
domains 4π, isotropic), sliced into S = 7 stencils at the sweep steps ∆T ∈ {5 × 10−3 , 10−2 , 1.5 ×
10−2 }. The unconditioned control pool is 18 (member, ∆T ) roots (FRC-{b2, b25, kf4, Re25k,
combo, 256}×3 steps); the conditioned run trained on an enlarged 41-root pool that adds further
forced members (FRC-{b0, b05, b075, b1}) and decaying-turbulence members. The       ... hygiene control
of Sec. 9.1 isolates that data difference on the 17 shared roots. Targets [Ṅ , N̈ , N ] are computed
once per window by the exact spectral chain rule (5) (forcing rebuilt per member from its manifest;
the targets themselves are forcing-free since m ≥ 1). The ∆T sweep deliberately includes near-wall
and past-wall tiers (Re25k at 1.5 × 10−2 exceeds its ∆T ⋆ ).

Pipeline safeguards (each one earned by a failure).
• By-window splits. Whole windows go to one split; anchored slices within a window are near-
  duplicates ∼ hf apart, so per-sample random splits leak and chronological splits shift covariates.
• Quiescent-window filter. Early quasi-zonal spin-up flow has J (ψ, ω) ≈ 0 and derivative targets
  104 –105 × smaller than developed flow; per-sample relative losses explode on them and the
  optimizer’s cheapest move is shrinking all predictions toward zero. Windows are dropped when
  the target norm falls below 10−2 × the member median or the stack roughness freezes. Spin-up
  duration is β-dependent (9–33% of windows across members), so per-member start times are
  set by a developed-flow criterion, and the filter is the safety net.
• Multigrid handling. Batches are shape-homogeneous but the dx-rescale of the stencils is strictly
  per-sample (equal shape does not imply equal dx); dealias masks are per-shape; anisotropic
  members are refused.
• Precision. All closure-pipeline compute is float64: the raw target is O(∆T 3 ) ≈ 10−9 , below
  float32 machine epsilon.
                   ...      (Disk storage of inputs is float32, upcast at load; this floors only the
  deepest channel N , which carries 2.6% of the closure.)

Optimization. Offline: AdamW, lr 5 × 10−5 , weight decay 10−4 , cosine schedule, 200 epochs,
batch 4, float64, physics init, seed 0. Rollout fine-tune (reference run, whose epoch-33 checkpoint
produced all a posteriori results): warm start from the offline conditioned model; strides 1, 2, 3
(∆T = stride × 5 × 10−3 ); unroll schedule 4:6, 8:8, 12:8, 16:8, 21:10 with per-stride clamps 21/7/3;


                                                   15
\end{Verbatim}
\newpage
% ---- Original PDF page 16 ----
\begin{Verbatim}[fontsize=\fontsize{6.35}{7.15}\selectfont,breaklines=false]
truncated BPTT window 4; lr 5 × 10−5 ; grad-clip 1.0; batch 1 with member-pure accumulation; 40
epochs.

9. Results.

9.1. A priori accuracy and the floors. Pooled test relative L2 per order (means; medians in
parentheses where they differ materially):

                                                                                                    ...
 model                                                pooled        Ṅ              N̈              N
 unconditioned control (deriv7_filtered_floor0.1)      0.269       0.058           0.186           0.563
 conditioned (cond_local, best ep63)                   0.214   0.074 (0.054)   0.138 (0.093)   0.430 (0.178)
 hygiene control (equal 17-root data, ep107)           0.250       0.059           0.178           0.511

Three readings. (a) The unconditioned N̈ plateau 0.186 matches the predicted pooled raw floor
0.175 of (38): the control is data-efficient but variance-limited, as derived. (b) Conditioning is real
on equal data: the conditioned model beats the hygiene control on N̈ on 16/17 shared roots (e.g. kf4
across the three ∆T : 0.058/0.065/0.057 vs 0.085/0.088/0.086), and the improvement concentrates
at low κ (the ε(κ) profile is U-shaped, not a high-k cliff). (c) The predicted conditioned ceilings were
not reached: kf4@ 1.5 × 10−2 plateaued at 0.065 vs the predicted 0.023, pooled 0.214 vs ≈ 0.05;
the run was stopped at plateau (epoch 78 of 300). The prediction is open, not falsified—the gap
is the difference between the information-theoretic floor of the conditioned family and what this
optimizer/parameterization extracted—but we report it as unmet.




         Figure placeholder: per-(member,∆T ,order) a priori rel-L2 heatmap, conditioned
                    vs control, with the per-tier Wiener floors overlaid. Source:
                                training/eval_deriv_by_root.py,
                       Theoretical_guarantees/check_wiener_floor.py.




9.2. A posteriori rollout: accuracy. Protocol: 16 coarse closed steps against fine-RK4 truth
(hf = 10−5 ; K = 500–1500), improvement = relLbare
                                              2    /relLclosed
                                                        2      at t = 16∆T ; 10 initial-condition
draws per ∆T across three members—kf4 and combo at 5122 , FRC-256 at 2562 —with combo held
out of the fine-tune entirely.

Analytic ceiling. The network-free R3 closure (exact chain-rule Ṅ , N̈ ; no learned component)
is stable in every arm and sets the honest upper bound: 71.4× / 35.1× / 132.6× over bare at
∆T = 5×10−3 /10−2 /1.5×10−2 . A one-step truncation-ratio computation predicts 66.6/30.7/19.1×:
the two agree within 7–14% at the two smaller steps, while the 132.6× at 1.5 × 10−2 exceeds its
one-step ceiling because the bare baseline there enters super-linear error growth—an accumulation
property of the baseline, not of the closure.

Learned closure.      After rollout fine-tuning, over the 10-draw replication:


                                                 16
\end{Verbatim}
\newpage
% ---- Original PDF page 17 ----
\begin{Verbatim}[fontsize=\fontsize{6.35}{7.15}\selectfont,breaklines=false]
          ∆T           stable   improvement over bare, median [min, max]   vs analytic ceiling
                −3
          5 × 10       10/10              15.96× [0.92, 23.57]                  71.4×
          10−2          9/10               0.50× [0.01, 6.17]                   35.1×
          1.5 × 10−2    6/10               0.65× [0.33, 1.89]              19.1× (one-step)

The 16× gain at the operating step replicates across draws and members, including the held-out
member (its worst draw is parity, 0.92). At the two larger steps the learned closure is at or
below parity where it is stable, and stability itself degrades to 60% of draws at 1.5 × 10−2 . For
calibration: the pre-fine-tune model blew up at the two larger steps (steps 7 and 13) and was stable
but below parity (0.72×) at the operating step—the fine-tune converted 0.72× into 16× there and
only partially cured the envelope edge. The gap between 16× and the 71× analytic ceiling is the
direct cost of the remaining N̈ estimation error (Prop. 7).




         Figure placeholder: rollout error growth, bare vs closed vs analytic-R3 , per ∆T ;
         vorticity fields and differences at t = 16∆T ; improvement-factor box plots over
                 the 10 draws. Source: diagnostics/rollout_aposteriori.py,
                          diagnostics/consolidate_apost_cases.py.




9.3. A posteriori rollout: stability, the horizon wall, and the mechanism. Three stability
findings, reported in full.

The instability is network-specific. The analytic R3 arm is stable at every (∆T, IC) tested,
at CFL up to 0.85; the learned arm can diverge even at CFL 0.56. Removing the N̈ channel from
the learned correction restores stability in every arm but also removes all gain (1.00× everywhere):
the destabilizer and the value live in the same channel.

Mechanism: history contamination. The blow-up is a feedback loop specific to the memory of
the method: network noise enters the correction, contaminates the stored 7-level history, is amplified
by the 1/∆T k time-FD weights into noisier derivative estimates, and returns as a noisier correction.
The blow-up horizon tracks (steps until injected noise ≈ truncation signal) + ∼4. Aliasing is not the
carrier (the learned arm gets worse in a fully dealiased world; the analytic arm is indifferent), and
the radial enstrophy-injection component is small and ∆T -independent—a dissipative projection
of the correction (P0 experiment) shifts blow-ups by ≤ ±1 step and fails its pre-registered stability
bars, ruling out the cheap inference-time fix. The within-shell (phase) component identified by the
certificate is the active target.

The horizon wall. Extending the 5 × 10−3 arm from 16 to 64 and 128 steps: 0/10 draws stable
in both arms, blow-ups at steps 21–38, i.e. at ≈ 2× the fine-tune’s free horizon. The replicated
16× gain is therefore horizon-limited: valid within the trained window, not unconditional. This
is the sharpest open problem; the certificate penalty (Sec. 7.4) attacks exactly this—a per-step



                                                  17
\end{Verbatim}
\newpage
% ---- Original PDF page 18 ----
\begin{Verbatim}[fontsize=\fontsize{6.35}{7.15}\selectfont,breaklines=false]
contraction certificate is horizon-free by construction—and the corresponding λ-sweep is running
at the time of writing.




         Figure placeholder: blow-up step vs horizon, 10 draws, 16/64/128 steps, with the
         free-horizon line; certificate |Geff (κ)| histograms before/after the penalty. Source:
                             diagnostics/aposteriori_stability.py,
                                 training/wiener_certificate.py.




9.4. Cost. Measured on the production stepper (5122 ): the closed step costs 8 FFTs vs the bare
5, plus ≈ 170 MACs/point of network—a per-step overhead of ≈ 60%, against the K-fold step-
count reduction of (24). At the replicated operating point (∆T = 5 × 10−3 , 16× error reduction
on-horizon), matching the closed accuracy with the bare scheme would require ≈ 4× smaller steps
(second-order scheme), so the realized speedup at matched accuracy is ≈ (4 × 5)/8 = 2.5× against
bare AB2–CN2 on the trained horizon; the ideal bound (24) is approached only as the learned
N̈ error approaches the analytic ceiling. The conditioned context costs zero additional FFTs at
rollout (spectral states reused); training-side costs are modest (rollout fine-tune ≈ 2–3 GPU-h;
the offline ensemble run ≈ 20 GPU-h; truth generation dominates and is cached across all experi-
ments). [TODO: tighten the realized-speedup accounting with wall-clock numbers from
rollout_timed_pareto.py on the final model.]

10. Conclusions and future work. We posed the time-discretization error of spectral AB2–CN2
as a closure problem and carried it from closed-form truncation theory to a deployed, certificate-
regularized learned scheme. The theory side is complete and verified: LTEs of the scheme family
to O(h6 ) (Props. 1–3), a finite and universal convergence radius ∆T ⋆ = 2.08 τλ (Prop. 5), an inner
stencil wall (Prop. 6), benign 1:1 error propagation (Prop. 7), the margin asymmetry that makes
fine-RK4 the right target (Sec. 5.2), and an exact Wiener analysis of the estimator that derives—
and prices—physics conditioning (Sec. 6.2). The practice side is honest: the analytic R3 closure
alone is unconditionally stable in our tests and worth up to two orders of magnitude; the learned
closure replicates 16× at the operating step on the trained horizon, beats its matched control on
the dominant N̈ channel on 16 of 17 shared roots as an estimator, but has not reached its predicted
conditioned floor, is at parity at the envelope edge, and remains horizon-limited by a diagnosed
history-contamination loop.

Future work. (i) Certificate-regularized training (λ-sweep in flight): convert the horizon-limited
gain into a per-step contraction guarantee. (ii) Conditioned-floor gap: close the factor 2–3 between
the trained conditioned model and its information-theoretic ceiling (longer training, richer σ̂ fea-
tures, angular harmonics for the r3 remainder). (iii) Weight-tied recursion cell: one cell unrolled p
times generalizes the estimator to any order N (p) with the exact ∂t recursion—no inner wall—at
the price of a ∆−1 per order; the exact-FFT vs learned-local ∆−1 decision is open, and is the
same decision as the deferred flux-form vs advective Jacobian consistency. (iv) Decaying-turbulence
members (non-stationary: the τλ identity needs detrending) and obstacle flows. (v) Generalization
across schemes: the entire construction used only (r, s) and the AB bracket; matching any linear
multistep method to its better-resolved self is mechanical from Prop. 3’s template.


                                                 18
\end{Verbatim}
\newpage
% ---- Original PDF page 19 ----
\begin{Verbatim}[fontsize=\fontsize{6.35}{7.15}\selectfont,breaklines=false]
Reproducibility. The        spectral      QG     solver      follows    Suresh   Babu    et    al.
[15].              Truncation       operators:           analysis/validate_ab2cn2_vs_truth.py;
magnitudes         and        radii:              analysis/measure_truncation_magnitudes.py,
Theoretical_guarantees/convergence_radius.py; microscale: check_decorrelation_time.py;
stencil     walls:           fd_depth_check.py,            temporal_fd_floor_deep.py;          er-
ror       propagation:               closure_error_propagation.py;              Wiener      floors:
check_wiener_floor.py;          stability    regions:          analysis/ab_stability_regions.py,
analysis/vn_stability.py,                 Theoretical_guarantees/imex_stability_regions.py;
certificate:    training/wiener_certificate.py;             training:     training/train_deriv.py,
training/train_deriv_rollout.py;             evaluation:          training/eval_deriv_by_root.py,
diagnostics/rollout_aposteriori.py. [TODO: code/data availability statement +
repository link.]

A. RK4 stage expansions. Throughout, f = G(ω n ), J = G′ (ω n ) = L + 2B(ω n , ·), and the exact
identity (7): G(ω n + p) = f + Jp + B(p, p), B symmetric bilinear.
                                         2
A.1. k2 . p = h2 f : k2 = f + h2 Jf + h4 B(f, f ).
                                                            2        3
A.2. k3 . p = h2 k2 . Linear part Jp = h2 Jf + h4 J 2 f + h8 JB(f, f ); quadratic part B(p, p) =
h2           h3            h4                 h4
4 B(f, f ) + 4 B(f, Jf ) + 8 B(f, B(f, f )) + 16 B(Jf, Jf ) + O(h ). Hence
                                                                 5

                    2               3                3             4                  4
 k3 = f + h2 Jf + h4 J 2 f +B(f, f ) + h8 JB(f, f )+ h4 B(f, Jf )+ h8 B(f, B(f, f ))+ h16 B(Jf, Jf ). (49)

A.3. k4 . p = hk3 . Collecting linear and quadratic parts through h4 ,
                                                                                       
    k4 = f + hJf + h2 12 J 2 f + B(f, f ) + h3 14 J 3 f + 41 JB(f, f ) + B(f, Jf )
                                                                                                   
          + h4 18 J 2 B(f, f ) + 41 JB(f, Jf ) + 12 B(f, J 2 f ) + 12 B(f, B(f, f )) + 14 B(Jf, Jf ) .   (50)

No stage builds the depth-four chain, so the h4 coefficient of J 4 f in k4 is zero; assembling h6 (k1 +
2k2 + 2k3 + k4 ) and subtracting from the exact expansion gives Prop. 1, per elementary differential:

                        elementary differential     exact ×h5    RK4 ×h5      τRK4 × h5

                        J 4f                         1/120           0         1/120
                          2
                        J B(f, f )                    1/60         1/48       −1/240
                        JB(f, Jf )                    1/20         1/24        1/120
                               2
                        B(f, J f )                    1/15         1/12        −1/60
                        B(f, B(f, f ))                2/15         1/8         1/120
                        B(Jf, Jf )                    1/20         1/16        −1/80


References.
 [1] U. M. Ascher, S. J. Ruuth, and B. T. R. Wetton. Implicit–explicit methods for time-dependent
     partial differential equations. SIAM J. Numer. Anal., 32(3):797–823, 1995.

 [2] Y. Bar-Sinai, S. Hoyer, J. Hickey, and M. P. Brenner. Learning data-driven discretizations for
     partial differential equations. Proc. Natl. Acad. Sci., 116(31):15344–15349, 2019.

 [3] J. P. Boyd. Chebyshev and Fourier Spectral Methods. 2nd ed., Dover, 2001.

                                                       19
\end{Verbatim}
\newpage
% ---- Original PDF page 20 ----
\begin{Verbatim}[fontsize=\fontsize{6.35}{7.15}\selectfont,breaklines=false]
 [4] J. C. Butcher. Numerical Methods for Ordinary Differential Equations. 3rd ed., Wiley, 2016.

 [5] C. Canuto, M. Y. Hussaini, A. Quarteroni, and T. A. Zang. Spectral Methods: Fundamentals
     in Single Domains. Springer, 2006.

 [6] A. Charous and P. F. J. Lermusiaux. Dynamically orthogonal Runge–Kutta schemes with
     perturbative retractions for the dynamical low-rank approximation. SIAM J. Sci. Comput.,
     45(2):A872–A897, 2023.

 [7] J. Frank, W. Hundsdorfer, and J. G. Verwer. On the stability of implicit-explicit linear multi-
     step methods. Appl. Numer. Math., 25(2–3):193–205, 1997.

 [8] Y. Guan, A. Chattopadhyay, A. Subel, and P. Hassanzadeh. Stable a posteriori LES of 2D
     turbulence using convolutional neural networks. J. Comput. Phys., 458:111090, 2022.

 [9] A. Gupta and P. F. J. Lermusiaux. Neural closure models for dynamical systems. Proc. R.
     Soc. A, 477:20201004, 2021.

[10] E. Hairer, S. P. Nørsett, and G. Wanner. Solving Ordinary Differential Equations I: Nonstiff
     Problems. 2nd ed., Springer, 1993.

[11] D. Kochkov, J. A. Smith, A. Alieva, Q. Wang, M. P. Brenner, and S. Hoyer. Machine learning–
     accelerated computational fluid dynamics. Proc. Natl. Acad. Sci., 118(21):e2101784118, 2021.

[12] J.-L. Lions, Y. Maday, and G. Turinici. Résolution d’EDP par un schéma en temps “pararéel”.
     C. R. Acad. Sci. Paris, 332:661–668, 2001.

[13] K. Srinivasan, M. D. Chekroun, and J. C. McWilliams. Turbulence closure with small, lo-
     cal neural networks: forced two-dimensional and β-plane flows. J. Adv. Model. Earth Syst.,
     16:e2023MS003795, 2024.

[14] H. J. Stetter. The defect correction principle and discretization methods. Numer. Math.,
     29:425–443, 1978.

[15] A. N. Suresh Babu, A. Sadam, and P. F. J. Lermusiaux. [TODO: title]. arXiv:2508.06678,
     2025.

[16] G. K. Vallis. Atmospheric and Oceanic Fluid Dynamics. 2nd ed., Cambridge University Press,
     2017.

[17] R. F. Warming and B. J. Hyett. The modified equation approach to the stability and accuracy
     analysis of finite-difference methods. J. Comput. Phys., 14:159–179, 1974.

[18] N. Wiener. Extrapolation, Interpolation, and Smoothing of Stationary Time Series. MIT Press,
     1949.




                                                20
\end{Verbatim}
\end{document}
